% !TEX spellcheck = en_US
% !TEX spellcheck = LaTeX
\documentclass[letterpaper,10pt,english]{article}
\usepackage{%
amsfonts,%
amsmath,%
amssymb,%
amsthm,%
babel,%
bbm,%
%biblatex,%
caption,%
centernot,%
color,%
enumerate,%
%enumitem,%
epsfig,%
epstopdf,%
etex,%
fancybox,%
framed,%
fullpage,%
%geometry,%
graphicx,%
hyperref,%
latexsym,%
mathptmx,%
mathtools,%
multicol,%
paralist,%
pgf,%
pgfplots,%
pgfplotstable,%
pgfpages,%
proof,%
psfrag,%
%subfigure,%
tcolorbox,%
tfrupee,%
tikz,%
times,%
%ulem,%
url,%
xcolor,%
mathpazo,%
}

\definecolor{shadecolor}{gray}{.95}%{rgb}{1,0,0}
\usepackage[margin=1in,top=0.75in]{geometry}
\usepackage[mathscr]{eucal}
\usepgflibrary{shapes}
\usepgfplotslibrary{fillbetween}
\usetikzlibrary{%
arrows,%
backgrounds,%
chains,%
decorations.pathmorphing,% /pgf/decoration/random steps | erste Graphik
decorations.text,% 
matrix,%
positioning,% wg. " of "
fit,%
patterns,%
petri,%
plotmarks,%
scopes,%
shadows,%
shapes.misc,% wg. rounded rectangle
shapes.arrows,%
shapes.callouts,%
shapes%
}

%\pgfplotsset{compat=newest} %<------ Here
\pgfplotsset{compat=1.11} %<------ Or use this one

\theoremstyle{plain}
\newtheorem{thm}{Theorem}[section]
\newtheorem{lem}[thm]{Lemma}
\newtheorem{prop}[thm]{Proposition}
\newtheorem{cor}[thm]{Corollary}
\newtheorem{clm}[thm]{Claim}

\theoremstyle{definition}
\newtheorem{defn}[thm]{Definition}
\newtheorem{conj}[thm]{Conjecture}
\newtheorem{exmp}[thm]{Example}
\newtheorem{axiom}[thm]{Axiom}
\newtheorem{assum}[thm]{Assumption}
\newtheorem{exerc}[thm]{Exercise}
\newtheorem*{defin}{Definition}

\theoremstyle{remark}
\newtheorem{rem}{Remark}
\newtheorem{note}{Note}

\newcommand{\iid}{\emph{i.i.d.}\ }
\newcommand{\Cov}{\operatorname{Cov}}
%\newcommand{\det}{\operatorname{det}}
%\newcommand{\Real}{\mathbb{R}}
\newcommand{\tr}{\operatorname{tr}}
\newcommand{\Var}{\operatorname{Var}}
\newcommand{\cov}{\operatorname{cov}}

%\renewcommand{\proof}[1]{\begin{proof}#1\end{proof}}
\newcommand{\EQ}[1]{\begin{equation*}#1\end{equation*}}
\newcommand{\EQN}[1]{\begin{equation}#1\end{equation}}
\newcommand{\eq}[1]{\begin{align*}#1\end{align*}}
\newcommand{\meq}[2]{\begin{xalignat*}{#1}#2\end{xalignat*}}
\newcommand{\norm}[1]{\left\lVert#1\right\rVert}
\newcommand{\inner}[1]{\left\langle #1\right\rangle}
\newcommand{\abs}[1]{\left\lvert#1\right\rvert}
\newcommand{\expect}[1]{\mathbb{E}\left[{#1}\right]}
\newcommand{\prob}[1]{\mathbb{P}\left[{#1}\right]}
\newcommand{\given}{\; \big\vert \;} 
\newcommand{\set}[1]{\left\{#1\right\}} 
\newcommand{\indicator}[1]{1_{\set{#1}}} 
\newcommand{\Ind}[1]{\mathbbm{1}_{#1}} 
\newcommand{\SetIn}[1]{\mathbbm{1}_{\set{#1}}} 
\newcommand{\red}[1]{\textcolor{red}{#1}} 
\newcommand{\floor}[1]{\left\lfloor#1\right\rfloor}
\newcommand{\ceil}[1]{\left\lceil#1\right\rceil}


\newcommand{\C}{\mathbb{C}}
\newcommand{\D}{\mathbb{D}}
\newcommand{\E}{\mathbb{E}}
\newcommand{\F}{\mathbb{F}}
\newcommand{\N}{\mathbb{N}}
\renewcommand{\P}{\mathbb{P}}
\newcommand{\Q}{\mathbb{Q}}
\newcommand{\R}{\mathbb{R}}
\newcommand{\Z}{\mathbb{Z}}

\newcommand{\bA}{\mathbf{A}}
\newcommand{\bI}{\mathbf{I}}
\newcommand{\bU}{\mathbf{1}}
\newcommand{\bx}{\mathbf{x}}
\newcommand{\bX}{\mathbf{X}}
\newcommand{\bZ}{\mathbf{Z}}


\newcommand{\cA}{\mathcal{A}}
\newcommand{\cB}{\mathcal{B}}
\newcommand{\cC}{\mathcal{C}}
\newcommand{\cF}{\mathcal{F}}
\newcommand{\cG}{\mathcal{G}}
\newcommand{\cM}{\mathcal{M}}
\newcommand{\cN}{\mathcal{N}}
\newcommand{\cP}{\mathcal{P}}
\newcommand{\cT}{\mathcal{T}}
\newcommand{\cX}{\mathcal{X}}
\newcommand{\cY}{\mathcal{Y}}

\newcommand{\sA}{\mathscr{A}}
\newcommand{\sB}{\mathscr{B}}
\newcommand{\sC}{\mathscr{C}}
\newcommand{\sE}{\mathscr{E}}
\newcommand{\sF}{\mathscr{F}}
\newcommand{\sG}{\mathscr{G}}
\newcommand{\sH}{\mathscr{H}}
\newcommand{\sL}{\mathscr{L}}
\newcommand{\sS}{\mathscr{S}}
\newcommand{\sT}{\mathscr{T}}
\newcommand{\sX}{\mathscr{X}}
\newcommand{\sY}{\mathscr{Y}}
\newcommand{\sZ}{\mathscr{Z}}

\newcommand{\tS}{\tilde{S}}
% Debug
\newcommand{\todo}[1]{\begin{color}{blue}{{\bf~[TODO:~#1]}}\end{color}}

% a few handy macros

\renewcommand{\le}{\leqslant}
\renewcommand{\ge}{\geqslant}
\newcommand\matlab{{\sc matlab}}
\newcommand{\goto}{\rightarrow}
\newcommand{\bigo}{{\mathcal O}}
%\newcommand{\half}{\frac{1}{2}}
%\newcommand\implies{\quad\Longrightarrow\quad}
\newcommand\reals{{{\rm l} \kern -.15em {\rm R} }}
\newcommand\complex{{\raisebox{.043ex}{\rule{0.07em}{1.56ex}} \hskip -.35em {\rm C}}}


% macros for matrices/vectors:

% matrix environment for vectors or matrices where elements are centered
\newenvironment{mat}{\left[\begin{array}{ccccccccccccccc}}{\end{array}\right]}
\newcommand\bcm{\begin{mat}}
\newcommand\ecm{\end{mat}}

% matrix environment for vectors or matrices where elements are right justifvied
\newenvironment{rmat}{\left[\begin{array}{rrrrrrrrrrrrr}}{\end{array}\right]}
\newcommand\brm{\begin{rmat}}
\newcommand\erm{\end{rmat}}

% for left brace and a set of choices
%\newenvironment{choices}{\left\{ \begin{array}{ll}}{\end{array}\right.}
\newcommand\when{&\text{if~}}
\newcommand\otherwise{&\text{otherwise}}
% sample usage:
%\delta_{ij} = \begin{choices} 1 \when i=j, \\ 0 \otherwise \end{choices}


% for labeling and referencing equations:
\newcommand{\eql}{\begin{equation}\label}
\newcommand{\eqn}[1]{(\ref{#1})}
% can then do
%\eql{eqnlabel}
%...
%\end{equation}
% and refer to it as equation \eqn{eqnlabel}.
% some useful macros for finite difference methods:
\newcommand\unp{U^{n+1}}
\newcommand\unm{U^{n-1}}
% for chemical reactions:
\newcommand{\react}[1]{\stackrel{K_{#1}}{\rightarrow}}
\newcommand{\reactb}[2]{\stackrel{K_{#1}}{~\stackrel{\rightleftharpoons}
 {\scriptstyle K_{#2}}}~}
\makeatletter
\def\th@plain{%
\thm@notefont{}% same as heading font
\itshape % body font
}
\def\th@definition{%
\thm@notefont{}% same as heading font
\normalfont % body font
}
\makeatother
\date{}


\title{Lecture-24: Poisson Point Processes}
\author{}

\begin{document}
\maketitle


%\section{Poisson and exponential random variables}
%\begin{defn}
%A non-negative integer valued random variable $N \in \Z_+$ is called \textbf{Poisson} if 
%for some constant $\lambda > 0$, we have 
%\EQ{
%P\{N = n\} = e^{-\lambda}\frac{\lambda^n}{n!}.
%}
%\end{defn}
%\begin{rem}
%It is easy to check that $\E N = \Var{N} = \lambda$. 
%Furthermore, the moment generating function $M_N_t = \E e^{t N} = e^{\lambda(e^{t} - 1)}$ exists for all $t \in \R$. 
%\end{rem}
%\subsection{Memoryless distribution}
%\begin{defn} 
%A random variable $X$ with continuous support on $\R_+$, is called \textbf{memoryless} if for all positive reals $t,s \in \R_+$, we have
%\EQ{
%  P\{X>s\} = P(\{ X > t+s\}\given \set{ X>t}). 
%}
%\end{defn}
%\begin{prop} The unique memoryless distribution function with continuous support on $\R_+$ is the exponential distribution.
%\end{prop}
%\begin{proof}
%Let $X$ be a random variable with a memoryless distribution function $F: \R_+ \to [0,1]$.  
%It follows that $\bar{F}(t) \triangleq 1 - F(t)$ satisfies the semi-group property 
%\EQ{
% \bar{F}(t+s) = \bar{F}(t)\bar{F}(s).
%}
%Since $\bar{F}(x) = P\{X > x\}$  is non-increasing in $x \in \R_+$, we have $\bar{F}(x) = e^{\theta x}$, for some $\theta < 0$ from Lemma~\ref{lem:semigroup}.
%\end{proof}

\section{Simple point processes} 
Consider the $d$-dimensional Euclidean space $\R^d$. 
The collection of Borel measurable subsets $\sB(\R^d)$ of the above Euclidean space is generated by sets $B(x) \triangleq \set{y \in \R^d: y_i \le x_i}$ for $x \in \R^d$.  
\begin{defn} 
A \textbf{simple point process} is a random countable collection of \emph{distinct points} $S: \Omega \to \sX^\N$, 
such that the distance $\norm{S_n} \to \infty$ as $n \to \infty$. 
\end{defn}
\begin{rem}
Since $S$ is a simple point process, each point $S_n$ is unique. 
Therefore, we can identify $S$ as a random set of points in $\sX$ and $S\cap A$ is the random set of points in $A$. 
\end{rem}
\begin{rem}
For any simple point process $S$, we have $P(\set{S_n = S_m \text{ for any } n \neq m}) = 0$ and $\abs{S\cap A}$ is finite almost surely for any bounded set $A \in \sB(\sX)$.
\end{rem}
%\begin{shaded}
%\begin{rem}
%Since $S$ is a simple point process, each point $S_n$ is unique. 
%Therefore, we can identify $S$ as a random set of points in $\sX$ and $S\cap A$ is the random set of points in $A$. 
%It follows that $\set{N(A) = n} = \set{\abs{S\cap A} = n}$. 
%\end{rem}
%\end{shaded}
\begin{shaded}
\begin{exmp}[Simple point process on the half-line] 
We can simplify this definition for $d=1$. 
When $\sX = \R_+$, one can order the points of the process $S: \Omega \to \R_+^\N$ to get ordered process $\tilde{S}: \Omega \to \R_+^\N$, 
such that $\tilde{S}_n = S_{(n)}$ is the $n$th order statistics of $S$. 
That is, $S_{(0)} \triangleq 0$, 
and 
$
S_{(n)} \triangleq \inf\set{S_k > S_{(n-1)}: k \in \N}. 
$ 
such that $S_{(1)} < S_{(2)} < \dots < S_{(n)} < \dots$, and $\lim_{n \in \N}S_{(n)} = \infty$. 
We will call this an arrival process. 
%The Borel measurable sets for $\R_+$ are generated by the collection of half-open intervals $\set{(0, t]: t \in \R_+}$. 
\end{exmp}
\end{shaded}
%\begin{shaded} 
%Point processes can model many interesting physical processes. 
%\begin{enumerate}
%\item Arrivals at classrooms, banks, hospital, supermarket, traffic intersections, airports etc.
%\item Location of nodes in a network, such as cellular networks, sensor networks, etc.
%\end{enumerate}
%\end{shaded}
\begin{defn}
Corresponding to a point process $S:\Omega\to\sX^\N$, 
we denote the number of points in a set $A \in \sB(\sX)$ by
\EQ{
N(A)\triangleq \abs{S\cap A} = \sum_{n \in \N}\Ind{A}(S_n), \text{ where we have }N(\emptyset) =0. 
}
%Then $N = (N(A): A \in \cB(\R^d))$ 
The resulting process $N: \Omega \to {\Z_+}^ {\sB(\sX)}$ is called a \textbf{counting process} for the point process $S: \Omega \to \sX^\N$. 
\end{defn} 
\begin{rem}
Let $A \in \sB(\sX)^k$ be a bounded partition of $B\in \sB(\sX)$. 
From the disjointness of $(A_1, \dots, A_k)$, we have 
\EQ{
N(B) = \sum_{n \in \N}\Ind{\cup_{i=1}^kA_i}(S_n)
= \sum_{n \in \N}\sum_{i=1}^k\Ind{A_i}(S_n)
= \sum_{i=1}^k\sum_{n \in\N}\Ind{A_i}(S_n)
= \sum_{i=1}^kN(A_i).
%= N(A_1) + \dots + N(A_k).
} 
%The result follows from the linearity of expectations, such that 
%$\Lambda(B) = \E N(B) = \sum_{i=1}^k\E N(A_i) = \sum_{i=1}^k\Lambda(A_i).$
\end{rem}
\begin{defn} 
A counting process is \textbf{simple} if the underlying point process is simple. %, and $N(\set{x}) \in \set{0,1}$ for all $x \in \R^d$. 
\end{defn}
\begin{rem}
For a simple counting process $N$, we have $N(\set{x}) \le 1$ almost surely for all $x \in \sX$.
\end{rem}
%\begin{shaded}
\begin{rem}
Let $N : \Omega \to {\Z_+}^{\sB(\sX)}$ be the counting process for the point process $S: \Omega \to \sX^\N$.  
\begin{enumerate}[i\_]
\item Note that the point process $S$ and the counting process $N$ carry the same information. 
\item The distribution of point process $S$ is completely characterized by the finite dimensional distributions of random vectors $(N(A_1), \dots, N(A_k))$ for any bounded sets $A_1, \dots, A_k \in \sB(\sX)$ and finite $k \in \N$.  
\end{enumerate}
\end{rem}
%\end{shaded}
\begin{shaded} 
\begin{exmp}[Simple point process on the half-line] 
Since the Borel measurable sets $\sB(\R_+)$ are generated by half-open intervals $\set{(0,t]: t \in \R_+}$, 
we denote the counting process by $N: \Omega \to {\Z_+}^{\R_+}$, 
where $N_t \triangleq N(0, t]= \sum_{n \in \N}\SetIn{S_n \in (0,t]}$ is the number of points in the half-open interval $(0,t]$. 
For $s < t$, the number of points in interval $(s,t]$ is $N(s,t] = N(0,t] - N(0,s] = N_t - N_s$. 
\end{exmp}
\end{shaded}
\begin{thm}[R\'{e}nyi] 
\label{thm:void}
Distribution of a simple point process $S:\Omega\to\sX^\N$ on a locally compact second countable space $\sX$ is completely determined by void probabilities $(P\set{N(A)=0}: A \in \sB(\sX))$. 
\end{thm}
\begin{proof}
It suffices to show that the finite dimensional distributions of $S$ on locally compact sets are characterized by void probabilities. 
\begin{enumerate}
\item[\bf{Step 1:}]  
We will show this by induction on the number of points $k$ in a bounded set $A \in \sB$. 
Let $A_1, \dots, A_k, B \in \sB(\sX)$ locally compact, then we will show that $u_k \triangleq P(\cap_{i=1}^k\set{N(A_i)> 0}\cap\set{N(B) = 0})$ can be computed from void probabilities.
From $k=1$, we have 
\EQ{
P\set{N(A_1) > 0, N(B) = 0} = P\set{N(B) = 0} - P\set{N(B\cup A_1) = 0}.
}
The induction can be proved by the recursive relation 
\EQ{
u_k
= P(\cap_{i=1}^{k-1}\set{N(A_i)> 0}\cap\set{N(B) = 0}) - P(\cap_{i=1}^{k-1}\set{N(A_i)> 0}\cap\set{N(A_k\cup B) = 0}).
} 
\item[\bf{Step 2:}]  
For any locally compact set $B \in \sB(\sX)$, there exists a sequence of nested partitions $B_n \triangleq (B_{n,j}: j \in [J_n])$ that eventually separates the points in $S\cap B$ as $n \to \infty$. 
We define the number of subsets of partition $(B_{n,j}: j \in [J_n])$ that consist of at least one point in $S\cap B$, as $H_n(B) \triangleq \sum_{j=1}^{J_n}\SetIn{N(B_{n,j}) > 0}$ where $H_n(B) \uparrow N(B)$ almost surely. 
\item[\bf{Step 3:}]  
We next show that for all locally compact sets $B_1, \dots, B_k \in \sB(\sX)$ and $j_1, \dots, j_k \in \N$, 
the probability $P(\cap_{i=1}^k\set{H_n(B_i) = j_i})$ can be expressed in terms of void probabilities. 
%There exists $n$ large enough so that we can find partitions $B^i_n$ for each $B_i$ such that each $B^i_{n,\ell}$ contains at maximum one point.  
We observe that 
\EQ{
P(\cap_{i=1}^k\set{H_n(B) = j_i})
= \sum_{ T_1, \dots, T_k \subseteq [J_n]: \abs{T_1} = j_1, \dots, \abs{T_k}=j_k}P\Big(\cap_{i=1}^k\cap_{j\in T_i}\set{N(B^i_{n,j})>  0}\cap\set{N(\cup_{j \notin \cup_{i=1}^kT_i}B^i_{n,j}) =0}\Big).
}
This can be expressed in terms of void probabilities by Step~1. 
\item[\bf{Step 4:}] For a simple point process, we have the following almost sure limit $\lim_n\cap_{i=1}^k\set{H_n(B_i) = j_i} = \cap_{i=1}^k\set{N(B_i) = j_i}$. 
The result follows from the continuity of probability.
\end{enumerate}
\end{proof}
\begin{rem} 
Recall that $\abs{A} = \int_{x \in A}dx$ is the volume of the set $A \in \sB(\R^d)$ and for any such $A$.
\end{rem}
\begin{defn}
The \textbf{intensity measure} $\Lambda:\sB(\sX)\to\R_+$ is defined for each bounded set $A \in \sX$ as its scaled volume 
in terms of the \textbf{intensity density} $\lambda: \R^d \to \R_+$, as 
\EQ{
\Lambda(A) \triangleq \int_{x \in A}\lambda(x) dx. 
} 
If the intensity density $\lambda(x) = \lambda$ for all $x \in \R^d$, then $\Lambda(A) = \lambda\abs{A}$. 
In particular for partition $A_1, \dots, A_k$ for a set $B$, we have $\Lambda(B) = \sum_{i=1}^k\Lambda(A_i)$. 
\end{defn}
%\begin{rem}
%From R\'{e}nyi's theorem, finite dimensional distribution of counting process of simple point processes is determined by void probabilities. 
%\end{rem}
%%%%%%%%%%%%%%%%%%%%%%%%%%%%%%%%%%%%%
\section{Poisson point process}
%%%%%%%%%%%%%%%%%%%%%%%%%%%%%%%%%%%%%
\begin{defn}
A non-negative integer valued random variable $N: \Omega \to \Z_+$ is called \textbf{Poisson} if 
for some constant $\lambda > 0$, we have 
\EQ{
P\set{N = n} = e^{-\lambda}\frac{\lambda^n}{n!}.
}
\end{defn}
\begin{rem}
It is easy to check that $\E N = \Var[N] = \lambda$. 
Furthermore, the moment generating function $M_{N_t} = \E e^{t N} = e^{\lambda(e^{t} - 1)}$ exists for all $t \in \R$. 
\end{rem}
\begin{cor}
\label{cor:ExpPoisson}
A simple counting process $N: \Omega \to \Z_+^{\sB(\sX)}$ has Poisson marginal distribution with intensity measure $\Lambda:\sB(\sX)\to\R_+$ if and only if void probabilities are exponential with the same intensity measure $\Lambda$. 
\end{cor}
\begin{proof}
It is clear that if the marginal distribution of the counting process $N$ is Poisson with intensity measure $\Lambda$, then the void probability $P\set{N(A) = 0} = e^{-\Lambda(A)}$ is exponential for any bounded set $A \in \sB(\sX)$. 

Conversely, we assume that the void probabilities are exponentially distributed with intensity measure $\Lambda$. 
It follows from the linearity of intensity measure that for any finite, bounded, and disjoint sets $B_1, \dots, B_k \in \sB(\sX)$, we have 
\EQ{
P(\cap_{i=1}^k\set{N(B_i)=0}) 
= P\set{N(\cup_{i=1}^kB_i) = 0}
= e^{-\Lambda(\cup_{i=1}^kB_i)} 
= \prod_{i=1}^ke^{-\Lambda(B_i)}
= \prod_{i=1}^kP\set{N(B_i) = 0}.
}
That is, the Bernoulli random vector $(\SetIn{N(B_i) = 0}: i \in [k])$ is independent for any finite $k \in \N$ and bounded disjoint $\sB(\sX)$ measurable sets $B_1, \dots, B_k$. 
Next we consider a set $B \in \sB(\sX)$ and a partition $B_n \triangleq (B_{n,j}: j \in [J_n])$ of $B$ such that $\Lambda(B_{n,j}) = \frac{\Lambda(B)}{J_n}$ for all $j \in [J_n]$. 
It follows that $H_n(B) \triangleq \sum_{j=1}^{J_n}\SetIn{N(B_{n,j})> 0}$ is the sum of $J_n$ \iid Bernoulli random variables with success probability $p_n \triangleq 1-e^{-\Lambda(B)/J_n}$, 
and hence has a Binomial distribution with parameters $(J_n, p_n)$.   
%where $p_n \approx \frac{\Lambda(B)}{J_n}$
Therefore, 
%\EQ{
%P\set{H_n(B) = m} 
%= \sum_{T\subseteq [J_n]: \abs{T} = m}\prod_{j \in T}(1-e^{-\Lambda(B_{n,j})})\prod_{j\notin T}e^{-\Lambda(B_{n,j})}
%= e^{-\Lambda(B)} \sum_{T\in [J_n]: \abs{T} = m}\prod_{j \in T}(e^{\Lambda(B_{n,j})}-1).
%}
\EQ{
P\set{H_n(B) = m} = \frac{e^{-\Lambda(B)}}{m!}\binom{J_n}{m}(e^{\Lambda(B)/J_n}p_n)^m
= e^{-\Lambda(B)}\frac{J_n!}{(J_n-m)!}\Big(e^{\Lambda(B)/J_n}-1\Big)^m.
}
Recall that $H_n(B) \uparrow N(B)$ as $n\to\infty$ in the proof of R\'{e}nyi's Theorem, 
and $\lim_{n\to\infty}J_n = \infty$ and $\lim_{n\in\N}\abs{B_{n,j}}= 0$. 
Thus, $\lim_{n\to\infty}\frac{J_n!}{(J_n-m)!}(e^{\Lambda(B)/J_n}-1)^m = \Lambda(B)^m$.
Taking limit $n\to\infty$ on both sides of the above equation, we get the result. 
%\EQ{
%P\set{N(B) = m} = e^{-\Lambda(B)}\sum_{T\in [J_n]: \abs{T} = m}\prod_{j \in T}\Lambda(B_{n,j})\prod_{j\notin T}(1-\Lambda(B_{n,j}))
%= e^{-\Lambda(B)}\frac{\Lambda(B)^m}{m!}.
%}
%We will skip the proof of converse that if void probabilities are exponential, then the marginal distribution is Poisson. 
%Since $A_k$ are very small, we get $e^{\Lambda(A_k)}-1 \approx \Lambda(A_k)$. 
%Therefore, we write that 
%\EQ{
%P\set{N(A)=1} 
%= \sum_{k \in \N}e^{-\Lambda(A)}(e^{\Lambda(A_k)}-1)
%\approx \Lambda(A)e^{-\Lambda(A)}. 
%}
%By inductive hypothesis, let $P\set{N(A) = n-1} = e^{-\Lambda(A)}\frac{\Lambda(A)^{n-1}}{(n-1)!}$ for all bounded $A \in \sB(\sX)$ and $n \in \N$. 
%Then, we can write 
%\EQ{
%P\set{N(A) = n} 
%= e^{-\Lambda(A)}\frac{\Lambda(A)^{n-1}}{(n-1)!}(1-\frac{1}{n} 
%+ \frac{1}{n}\sum_{k \in \N}\left(e^{\Lambda(A_k)}\Big(1-\frac{\Lambda(A_k)}{\Lambda(A)}\Big)^{n-1}-1\right). 
%}
%From the fact that $\sum_{k \in \N}\Lambda(A_k) = \Lambda(A)$ and the approximation $e^{\Lambda(A_k)}\Big(1-\frac{\Lambda(A_k)}{\Lambda(A)}\Big)^{n-1}-1 \approx \Lambda(A_k) - (n-1)\frac{\Lambda(A_k)}{\Lambda(A)}$ for sufficiently small $A_k$, 
%we get the induction step  
%$
%P\set{N(A) = n} 
%\approx e^{-\Lambda(A)}\frac{\Lambda(A)^n}{n!}.
%$
\end{proof}
\begin{defn} 
A counting process $N:\Omega\to\Z_+^{\sB(\sX)}$ has the \textbf{completely independence property},  
if for any collection of finite disjoint and bounded sets $A_1, \dots, A_k \in \sB(\sX)$, 
the vector $(N(A_1), \dots, N(A_k)):\Omega \to \Z_+^k$ is independent. 
That is,
\EQ{
P\left(\bigcap_{i=1}^k\set{N(A_i) = n_i}\right) = \prod_{i=1}^kP\set{N(A_i) = n_i},\quad n \in \Z_+^k.
}
\end{defn}

\begin{defn}
A simple point process $S: \Omega \to \sX^\N$ is \textbf{Poisson point process}, 
if the associated counting process $N: \Omega \to \Z_+^{\sB(\sX)}$ has complete independence property and the marginal distributions are Poisson. 
\end{defn}
\begin{defn}
The \textbf{intensity measure} $\Lambda: \sB(\sX) \to \R_+$ of Poisson process $S$ is defined by $\Lambda(A) \triangleq \E N(A)$ for all bounded $A \in \sB(\sX)$. 
\end{defn}
\begin{rem}
Recall that for any partition $A \in \sB(\sX)^k$ of a bounded set $B \in \sB(\sX)$, we have $N(B) = \sum_{i=1}^kN(A_i)$ and therefore it follows from the linearity of expectations that $\Lambda(B) = \E N(B) = \sum_{i=1}^k\E N(A_i) = \sum_{i=1}^k\Lambda(A_i).$ 
Thus, this is a valid intensity measure. 
\end{rem}
\begin{rem}
For a Poisson process with intensity measure $\Lambda$, 
it follows from the definition that for any finite $k \in \Z_+$, 
and bounded mutually disjoint sets $A_1, \dots, A_k \in \sB(\sX)$, we have  
\EQ{
P\left(\cap_{i=1}^k\set{N(A_i) = n_i}\right) = \prod_{i=1}^k\left(e^{-\Lambda(A_i)}\frac{\Lambda(A_i)^{n_i}}{n_i!}\right),\quad n \in \Z_+^k.
} 
\end{rem} 
\begin{defn}
If the intensity measure $\Lambda$ of a Poisson process $S$ satisfies $\Lambda(A) = \lambda\abs{A}$ for all bounded $A \in \sB(\sX)$, then we call $S$ a \textbf{homogeneous Poisson point process} and $\lambda$ is its intensity.  
\end{defn} 
%\begin{defn} 
%A counting process $N:\Omega\to\Z_+^{\sB(\sX)}$ has the \textbf{completely independence property},  
%if for any collection of finite disjoint and bounded sets $A_1, \dots, A_k \in \sB(\sX)$, 
%the vector $(N(A_1), \dots, N(A_k)):\Omega \to \Z_+^k$ is independent. 
%That is,
%\EQ{
%P\bigcap_{i=1}^k\set{N(A_i) = n_i} = \prod_{i=1}^kP\set{N(A_i) = n_i},\quad n \in \Z_+^k.
%}
%\end{defn}
%\begin{rem}
%\end{rem}
%\begin{shaded}
%\begin{thm}[Poisson process on the half-line]
%A random process $N: \Omega \to {\R_+}^ {\Z_+}$ indexed by time $t$ is a Poisson process iff
%\begin{enumerate}[(a)]
%\item Starting with $N(0) = 0$, the process $N_t$ takes a non-negative integer value for all $t \in \R_+$;
%\item the increment $N(t + s) - N_t$ is surely nonnegative for any $s \in \R_+$;
%\item the increments $N(t_1), N(t_2) - N(t_1),  \dots , N(t_n) - N(t_{n-1})$ are \textit{independent} for any 
%$0 < t_1 < t_2 < \dots < t_{n-1} < t_n$; 
%\item the increment $N(t+s) - N_t$ is distributed as Poisson random variable with parameter $\Lambda((t, t+s])$. 
%\end{enumerate}
%The Poisson process is homogeneous with intensity $\lambda$, iff in addition to conditions $(a), (b), (c)$, 
%the distribution of the increment $N(t + s) - N_t$ depends on the value $s \in \R_+$ but is independent of $t \in \R_+$. 
%That, is the increments are \textit{stationary}. 
%\end{thm}
%\begin{proof} 
%We have already seen that definition of Poisson processes implies all four conditions. 
%We will show in a later lecture that the above four conditions imply that $N: \Omega \to {\Z_+}^{\R_+}$ is a simple Poisson counting process. 
%\end{proof}
%\end{shaded}
%%%%%%%%%%%%%%%%%%%%%%%%%%%%%%%%%%%%%
\section{Equivalent characterizations}
%%%%%%%%%%%%%%%%%%%%%%%%%%%%%%%%%%%%%
%\begin{shaded}
%\begin{rem}
%Let $\sX = \R^d$, then the process $S: \Omega \to \sX^\N$ is a Poisson point process iff 
%\begin{enumerate}[i\_]
%\item the counting process $N$ has complete independence property, and 
%\item for each bounded set $A \in \sB$, the random variable $N(A)$ is Poisson with parameter $\Lambda(A)$. 
%%\item for all $k \in \N$ and all bounded, mutually disjoint $A_1, \dots, A_k \in \sB$, the random vector $(N(A_1), \dots, N(A_k))$ is a vector of independent random variables. 
%\end{enumerate}
%%\item with parameters $\Lambda(A_1), \dots, \Lambda(A_k)$ respectively. 
%In particular, we have $\E N(A) = \Lambda(A)$ for all subsets $A \in \sB$. 
%\end{rem}
%\end{shaded}
\begin{thm}[Equivalences] 
\label{thm:equiv}
Following are equivalent for a simple counting process $N: \Omega \to {\Z_+}^{\sB(\sX)}$. 
\begin{enumerate}[i\_]
\item Process $N$ is Poisson with locally finite intensity measure $\Lambda$. 
\item For each bounded $A \in \sB(\sX)$, we have $P\set{N(A) = 0} = e^{-\Lambda(A)}$. 
\item %The intensity measure $\Lambda$ is bounded for bounded $A \in \sB$,  and 
For each bounded $A \in \sB(\sX)$,  the number of points $N(A)$ is a Poisson with parameter $\Lambda(A)$. 
\item Process $N$ has the completely independence property, and $\E N(A) = \Lambda(A)$ for all bounded sets $A \in \sB(\sX)$. 
\end{enumerate}
\end{thm}
\begin{proof}
We will show that $i\_ \implies ii\_ \implies iii\_ \implies iv\_ \implies i\_$. 
\begin{enumerate}
	\item[$i \implies ii\_$] It follows from the definition of Poisson point processes and definition of Poisson random variables.  
	\item[$ii \implies iii\_$] 
	From Corollary~\ref{cor:ExpPoisson}, we know that if void probabilities are exponential, then the marginal distributions are Poisson. 
	%From Theorem~\ref{thm:void} we know that void probabilities determine the entire distribution. 
%Further, we observe that 
%\EQ{
%\sum_{k\in \N}P\set{N(A) = k} = e^{-\Lambda(A)}\sum_{k\in \N}\frac{\Lambda(A)^k}{k!}.
%}
%For any finite $k$ and bounded disjoint sets $B_1, \dots, B_k \in \sB(\sX))$, 
%we have 
%\EQ{
%P(\cap_{i=1}^k\set{N(B_i)=0}) 
%= e^{-\Lambda(\cup_{i=1}^kB_i)} 
%= \prod_{i=1}^ke^{-\Lambda(B_i)}
%= \prod_{i=1}^kP\set{N(B_i) = 0}.
%}
%We further observe that $H_n(B) = \sum_{j=1}^{J_n}\SetIn{N(B_{n,j})> 0}$ is the sum of $J_n$ independent Bernoulli random variables with success probability $1-e^{-\Lambda(B_{n,j})}$. 
%Therefore, 
%\EQ{
%P\set{H_n(B) = m} 
%= \sum_{T\in [J_n]: \abs{T} = m}\prod_{j \in T}(1-e^{-\Lambda(B_{n,j})})\prod_{j\notin T}e^{-\Lambda(B_{n,j})}.
%%= e^{-\Lambda(B)} \sum_{T\in [J_n]: \abs{T} = m}\prod_{j \in T}(e^{\Lambda(B_{n,j})}-1).
%}
%Recall that $H_n(B) \uparrow N(B)$ as $n\to\infty$ in the proof of R\'{e}nyi's Theorem, 
%and $\lim_{n\in\N}\abs{B_{n,j}}= 0$. 
%Taking limit $n\to\infty$ on both sides of the above equation, we get
%\EQ{
%P\set{N(B) = m} = e^{-\Lambda(B)}\sum_{T\in [J_n]: \abs{T} = m}\prod_{j \in T}\Lambda(B_{n,j})
%= e^{-\Lambda(B)}\frac{\Lambda(B)^m}{m!}.
%}
\item[$iii \implies iv\_$] We will show this in two steps. 
\begin{itemize}
\item[Mean:] Since the distribution of random variable $N(A)$ is Poisson, it has mean $\E N(A) = \Lambda(A)$. 
\item[CIP: ] Consider a partition $A \in \sB^k$ for a bounded set $B\in\sB(\sX)$, then 
%$N(B) = N(A_1) + \dots + N(A_k)$. %Taking expectations on both sides, and from the linearity of expectation, we get 
$\Lambda(B) = \Lambda(A_1) + \dots + \Lambda(A_k)$. 
Consider all partitions $n\in\Z_+^k$ of a non-negative integer $m \in \Z_+$, to write 
\EQ{
P\set{N(B)=m} = \sum_{n_1 + \dots + n_k = m}P\set{N(A_1) = n_1, \dots, N(A_k) = n_k}. 
}
Using the definition of Poisson distribution, we can write the LHS of the above equation as 
\EQ{
P\set{N(B) = m} 
= e^{-\Lambda(B)} \frac{\Lambda(B)^m}{m!} 
= \prod_{i=1}^ke^{-\Lambda(A_i)} \frac{(\sum_{i=1}^k\Lambda(A_i))^m}{m!}.
}
Since the expansion of $(a_1 + \dots + a_k)^m = \sum_{n_1 + \dots+n_k = m}\binom{m}{n_1, \dots, n_k}\prod_{i=1}^ka_i^{n_i}$, 
we get 
\EQ{
P\set{N(B) = m} 
= \frac{1}{m!}\sum_{n_1 + \dots + n_ k =n}\binom{m}{n_1, \dots, n_k}\prod_{i=1}^ke^{-\Lambda(A_i)} \Lambda(A_i))^{n_i}
= \sum_{n_1+\dots+n_k = m}\left(\prod_{i=1}^ke^{-\Lambda(A_i)} \frac{\Lambda(A_i))^{n_i}}{n_i!}\right). 
}
Equating each term in the summation, we get $P\set{N(A_1) = n_1, \dots, N(A_k) = n_k} = \prod_{i=1}^kP\set{N(A_i) = n_i}$. 
\end{itemize}

\item[$iv \implies i\_$] 
From Corollary~\ref{cor:ExpPoisson}, if the void probability is exponential with intensity measure $\Lambda$, then the marginal distribution if Poisson with the same intensity measure. 
%Therefore, it suffices to show that $P\set{N(A) = 0} = e^{-\Lambda(A)}$ for all bounded $A \in \sB$. 
%To show this, 
We define $f:\sB(\sX)\to(-\infty, 0]$ by $f(A) \triangleq \ln P\set{N(A) = 0}$ for all bounded $A \in \sB(\sX)$. 
Then, we observe that for any partition $(A_1,\dots, A_k)$ of $A$, we have $f(\cup_{i=1}^kA_i) = \ln P\set{N(A)= 0} = \ln\prod_{i=1}^kP\set{N(A_i) = 0} = \sum_{i=1}^kf(A_i)$.
It follows that $-f:\sB(\sX)\to\R_+$ is an intensity measure, and $P\set{N(A)=0} = e^{f(A)}$. 
Since $\E N(A) = -f(A) = \Lambda(A)$, the result follows. 
\end{enumerate}
\end{proof}

%\red{
%\begin{thm} 
%Following are equivalent for a simple counting process $N$. 
%\begin{enumerate}[i\_]
%\item Process $N$ is Poisson with intensity measure $\Lambda$. 
%\item Process $N$ has the completely independence property, and $\E N(A) = \Lambda(A)$. 
%\item For any bounded $A \in \sB$, the intensity measure $\Lambda$ is bounded and $N(A)$ is Poisson with parameter $\Lambda(A)$. 
%\end{enumerate}
%\end{thm}
%\begin{proof} 
%We will show that $i\_ \implies ii\_ \implies iii\_ \implies I\_$. 
%\begin{enumerate}[i\_]
%\item[$i\_ \implies ii\_$] Let the process $N$ be Poisson with intensity measure $\Lambda$. 
%Then the complete independence property follows from the definition, 
%and since $N(A)$ is a Poisson distributed random variable with mean $\Lambda(A)$. 
%\item[$iI\_ \implies iii\_$] Consider a process $N$ with the completely independence property, and $\E N(A) = \Lambda(A)$. 
%Let $A \in \sB$ be bounded, then it follows that $\Lambda(A)$ is bounded. 
%From the fact that $\Lambda(A) = \sum_{i=1}^k\Lambda(A_i)$ for any partition $(A_1, \dots, A_k)$ of the set $A$, 
%we can write $\set{N(A) = 0} = \cap_{i=1}^k\set{N(A_i) = 0}$. 
%Therefore, from the complete independence property, we have 
%\EQ{
%P\set{N(A) = 0} = \prod_{i=1}^kP\set{N(A_i) = 0}. 
%}
%This is the semi-group property, from which it follows that 
%\meq{2}{
%&P\set{N(A) = 0} = e^{-\Lambda(A)},&&\sum_{k\in \N}P\set{N(A) = k} = e^{-\Lambda(A)}\sum_{k\in \N}\frac{\Lambda(A)^k}{k!}.
%}
%From which it follows that $N(A)$ is Poisson with parameter $\Lambda(A)$. 
%\item[$iii\_ \implies i\_$] Let $N$ be a simple counting process and $\Lambda$ an intensity measure such that for any $A \in \sB$, 
%$\Lambda(A)$ is bounded and $N(A)$ is Poisson with parameter $\Lambda(A)$. 
%This implies that for any partition $(A_1, \dots, A_k)$ of the set $A_i$, we have 
%\EQ{
%P\set{N(A) = 0} = \prod_{i=1}^kP\set{N(A_i) = 0}. 
%}
%Hence, we can show that the counting process $N$ has the complete independence property. 
%\end{enumerate}
%\end{proof}
%}
\begin{cor}[Poisson process on the half-line]
A random process $N: \Omega \to \Z_+^{\R_+}$ indexed by time $t \in \Z_+$ is the counting process associated with a one-dimensional Poisson process $S: \Omega \to \R_+^\N$ having intensity measure $\Lambda$ iff
\begin{enumerate}[(a)]
\item Starting with $N_0 = 0$, the process $N_t$ takes a non-negative integer value for all $t \in \R_+$;
\item the increment $N_s - N_t$ is surely nonnegative for any $s \ge t$;
\item the increments $N_{t_1}, N_{t_2} - N_{t_1},  \dots , N_{t_n} - N_{t_{n-1}}$ are \textit{independent} for any 
$0 < t_1 < t_2 < \dots < t_{n-1} < t_n$; 
\item the increment $N_s - N_t$ is distributed as Poisson random variable with parameter $\Lambda(t, s]$ for $s \ge t$.  
\end{enumerate}
The Poisson process is homogeneous with intensity $\lambda$, iff in addition to conditions $(a), (b), (c)$, 
the distribution of the increment $N_{t+s} - N_t$ depends on the value $s \in \R_+$ but is independent of $t \in \R_+$. 
That, is the increments are \textit{stationary}. 
\end{cor}
\begin{proof} 
We have already seen that definition of Poisson processes implies all four conditions. 
Conditions $(a)$ and $(b)$ imply that $N$ is a simple counting process on the half-line, 
condition $(c)$ is the complete independence property of the point process, 
and condition $(d)$ provides the intensity measure. 
The result follows from the equivalence $iv\_$ in Theorem~\ref{thm:equiv}. 
\end{proof}


%\appendix
%
%\section{Memoryless distribution}
%\begin{defn} 
%A random variable $X$ with continuous support on $\R_+$, is called \textbf{memoryless} if for all positive reals $t,s \in \R_+$, we have
%\EQ{
  %P\set{X>s} = P(\set{ X > t+s}\given \set{ X>t}). 
%}
%\end{defn}
%\begin{prop} The unique memoryless distribution function with continuous support on $\R_+$ is the exponential distribution.
%\end{prop}
%\begin{proof}
%Let $X$ be a random variable with a memoryless distribution function $F: \R_+ \to [0,1]$.  
%It follows that $\bar{F}(t) \triangleq 1 - F(t)$ satisfies the semi-group property 
%\EQ{
 %\bar{F}(t+s) = \bar{F}(t)\bar{F}(s).
%}
%Since $\bar{F}(x) = P\set{X > x}$  is non-increasing in $x \in \R_+$, we have $\bar{F}(x) = e^{\theta x}$, for some $\theta < 0$ from Lemma~\ref{lem:semigroup}.
%\end{proof}
%
%%\section{Functions with semigroup property}
%\begin{lem} 
%\label{lem:semigroup}
%A unique non-negative right continuous function $f: \R_+ \to \R_+$ satisfying the semigroup property
%\EQ{
 %f(t+s) = f(t)f(s), \text{ for all } t,s \in \R_+
 %}
 %is $f(t) = e^{\theta t}$, where $\theta = \log f(1)$.
%\end{lem}
%\begin{proof}
%Clearly, we have $f(0) = f^2(0)$. Since $f$ is non-negative, it means $f(0) = 1$. 
%By definition of $\theta$ and induction for $m,n \in \Z^+$, we see that 
%\begin{xalignat*}{3}
%&f(m) = f(1)^m = e^{\theta m},&& e^{\theta} = f(1) = f(1/n)^n.
 %\end{xalignat*}
%Let $q \in \Q$, then it can be written as $m/n, n \neq 0$ for some $m,n \in \Z^+$. 
%Hence, it is clear that for all $q \in \Q^+$, we have $f(q) = e^{\theta q}$.
%either unity or zero. Note, that $f$ is a right continuous function and is non-negative. 
%Now, we can show that $f$ is exponential for any real positive $t$ by taking a sequence of rational numbers $(q_n: n \in \N)$ decreasing to $t$. 
%From right continuity of $f$, we obtain 
%\EQ{
%f(t) = \lim_{q_n \downarrow t} f(q_n) =   \lim_{q_n \downarrow t} e^{\theta q_{n}}= e^{\theta t}.
%}
%\end{proof}
%%
%%\begin{lem} 
%%\label{lem:semigroup}
%%A unique non-negative right continuous function $f: \R_+ \times \R_+\to \R_+$ satisfying the property
%%\EQ{
%%f(0,t+s) = f(0,t)f(t,t+s), \text{ for all } t,s \in \R_+
%%}
%%is $f(t,t+s) = e^{\Lambda(t,t+s)}$. 
%%\end{lem}
%%\begin{proof}
%%Clearly, we have $f(0,0) = f^2(0,0)$ and $f(0,t) = f(0,0)f(0,t)$. 
%%Since $f$ is non-negative, it means $f(0,0) = 1$.  
%%By definition of $\theta$ and induction for $m,n \in \Z^+$, we see that 
%%\begin{xalignat*}{2}
%%&f(0,m) 
%%= \prod_{i=1}^mf(i-1,i) 
%%= e^{\theta m},&
%%&f(0,1) = \prod_{i=1}^nf(\frac{i-1}{n}, \frac{i}{n}).
%%\end{xalignat*}
%%Let $q \in \Q$, then it can be written as $m/n, n \neq 0$ for some $m,n \in \Z^+$. 
%%Hence, it is clear that for all $q \in \Q^+$, we have $f(q) = e^{\theta q}$.
%%either unity or zero. Note, that $f$ is a right continuous function and is non-negative. 
%%Now, we can show that $f$ is exponential for any real positive $t$ by taking a sequence of rational numbers $(q_n: n \in \N)$ decreasing to $t$. 
%%From right continuity of $f$, we obtain 
%%\EQ{
%%f(t) = \lim_{q_n \downarrow t} f(q_n) =   \lim_{q_n \downarrow t} e^{\theta q_{n}}= e^{\theta t}.
%%}
%%\end{proof}
%

\end{document}